\documentclass[12pt]{article}
\usepackage[margin=0.8in]{geometry}
\usepackage{amsmath}
\usepackage{xcolor}
\usepackage{tcolorbox}
\usepackage{enumitem}

\definecolor{highlight}{RGB}{255,230,153}

\title{\textbf{Slutsky Decomposition: The Simple Version}}
\author{}
\date{}

\begin{document}
\maketitle

\section*{Forget the Jargon. Here's What's Actually Happening:}

You buy coffee and donuts every week. Suddenly, coffee gets more expensive.

Obviously, you're going to buy less coffee. But \textbf{WHY} are you buying less? There are actually \textbf{TWO separate reasons}:

\begin{tcolorbox}[colback=blue!5!white,colframe=blue!75!black]
\textbf{Reason 1: "Coffee is now expensive compared to donuts"}

You think: "Wow, coffee costs way more than donuts now. I should buy more donuts and less coffee because I get more bang for my buck with donuts."

\vspace{0.2cm}
$\rightarrow$ This is the \textbf{SUBSTITUTION EFFECT}
\end{tcolorbox}

\begin{tcolorbox}[colback=red!5!white,colframe=red!75!black]
\textbf{Reason 2: "I'm poorer now"}

You think: "Damn, my paycheck is the same but everything feels more expensive. I need to cut back on everything."

\vspace{0.2cm}
$\rightarrow$ This is the \textbf{INCOME EFFECT}
\end{tcolorbox}

That's it. That's the whole thing.

\section*{Concrete Example: You and Your Coffee}

\subsection*{Initial Situation}

\begin{itemize}[leftmargin=*]
\item You have \$20/week to spend
\item Coffee costs \$2 per cup
\item Donuts cost \$2 each
\item You buy: \textbf{5 coffees + 5 donuts} (spending all \$20)
\item Your happiness level: Let's call it \textbf{100 happy points}
\end{itemize}

\subsection*{What Happens: Coffee Price Doubles to \$4}

You end up buying: \textbf{2 coffees + 6 donuts}

So you went from 5 coffees down to 2 coffees. That's a change of \textbf{-3 coffees}.

But let's break down WHY you cut back by 3 coffees...

\newpage

\subsection*{Breaking It Down}

\textbf{STEP 1: The Substitution Effect}

Imagine I said: "Hey, I know coffee just got expensive. Let me give you extra money so you can still afford to be exactly as happy as before (100 happy points)."

With this extra compensation, you'd have enough money to stay at 100 happy points. But even with this extra cash, coffee is still \textit{relatively expensive} compared to donuts now.

So you'd think: "Coffee is twice as expensive as donuts now (\$4 vs \$2). Even though I can afford to be just as happy, I should substitute toward donuts and away from coffee."

You'd buy: \textbf{4 coffees + 7 donuts} (with the extra money I gave you)

\vspace{0.3cm}

\colorbox{highlight}{\textbf{Substitution Effect = 4 - 5 = -1 coffee}}

\vspace{0.3cm}

You cut back by 1 coffee just because of the \textit{relative price change}, even though you had enough money to stay equally happy.

\vspace{0.5cm}

\textbf{STEP 2: The Income Effect}

Now I take away the extra money I gave you. You're back to your real \$20/week budget.

But remember, coffee is now \$4 instead of \$2. Your \$20 doesn't go as far anymore. You're effectively \textbf{poorer}.

Being poorer means you need to cut back on \textit{everything} (both coffee and donuts).

You go from: \textbf{4 coffees} (with compensation) down to \textbf{2 coffees} (without compensation)

\vspace{0.3cm}

\colorbox{highlight}{\textbf{Income Effect = 2 - 4 = -2 coffees}}

\vspace{0.3cm}

You cut back by another 2 coffees just because you're effectively poorer now.

\vspace{0.5cm}

\textbf{STEP 3: Add Them Up}

\begin{tcolorbox}[colback=green!10!white,colframe=green!75!black,title=The Slutsky Decomposition]
\textbf{Total change in coffee} = -3 coffees

\vspace{0.2cm}

This breaks down as:

\textbf{Substitution Effect:} -1 coffee (coffee is expensive relative to donuts)

\textbf{Income Effect:} -2 coffees (you're poorer overall)

\vspace{0.2cm}

Check: $-1 + (-2) = -3$ \checkmark
\end{tcolorbox}

\newpage

\section*{The Visual Version}

Think of it like this:

\vspace{0.3cm}

\begin{center}
\fbox{\begin{minipage}{0.9\textwidth}
\textbf{Original:} 5 coffees, 5 donuts, 100 happy points

\vspace{0.3cm}

$\downarrow$ \textit{Coffee price doubles, but I give you extra money to stay at 100 happy points}

\vspace{0.3cm}

\textbf{Compensated:} 4 coffees, 7 donuts, 100 happy points

\colorbox{yellow}{This move is the SUBSTITUTION EFFECT: -1 coffee}

\vspace{0.3cm}

$\downarrow$ \textit{I take away the extra money, back to real \$20 budget}

\vspace{0.3cm}

\textbf{Final:} 2 coffees, 6 donuts, maybe 60 happy points

\colorbox{yellow}{This move is the INCOME EFFECT: -2 coffees}
\end{minipage}}
\end{center}

\section*{Key Insights}

\begin{enumerate}[leftmargin=*]
\item \textbf{Substitution effect is ALWAYS negative} (or zero)
\begin{itemize}
\item When price goes up, you always substitute away from that good
\item It's about \textit{relative prices}
\item "Coffee vs donuts: which gives more happiness per dollar?"
\end{itemize}

\item \textbf{Income effect depends on the type of good}
\begin{itemize}
\item \textbf{Normal good} (like coffee): When you're poorer, you buy less
\begin{itemize}
\item Income effect is negative
\item Both effects work together → demand definitely falls
\end{itemize}
\item \textbf{Inferior good} (like instant ramen): When you're poorer, you buy MORE
\begin{itemize}
\item Income effect is positive
\item The two effects fight each other
\end{itemize}
\end{itemize}

\item \textbf{Special cases} (you'll see these on the exam):
\begin{itemize}
\item \textbf{Perfect complements} (like left shoes and right shoes): Substitution effect = 0
\begin{itemize}
\item You can't substitute! You always need them in fixed proportions
\item Entire change is income effect
\end{itemize}
\item \textbf{Quasilinear utility}: Income effect = 0 for the "special" good
\begin{itemize}
\item Entire change is substitution effect
\end{itemize}
\end{itemize}
\end{enumerate}

\newpage

\section*{How to Solve These Problems on the Exam}

\subsection*{The Recipe (Simple Version)}

\textbf{Given:} Price of good 1 rises from $p_1$ to $p_1'$

\begin{enumerate}[leftmargin=*]
\item \textbf{Find the original bundle}
\begin{itemize}
\item Use the standard optimal choice formulas
\item For Cobb-Douglas $u = x_1^a x_2^b$: spend $\frac{a}{a+b}$ of income on good 1
\item Calculate the utility level at this bundle
\end{itemize}

\item \textbf{Find the compensated bundle}
\begin{itemize}
\item Stay at the same utility level
\item But use the NEW price ratio
\item Mathematically: solve $MRS = \frac{p_1'}{p_2}$ while keeping $u = \bar{u}$
\end{itemize}

\item \textbf{Find the new bundle}
\begin{itemize}
\item Use the NEW prices
\item Use the ORIGINAL income
\item Just solve the regular problem again
\end{itemize}

\item \textbf{Calculate the effects}
\begin{itemize}
\item Sub effect = compensated - original
\item Income effect = new - compensated
\item Total = new - original (should equal sub + income)
\end{itemize}
\end{enumerate}

\subsection*{Cobb-Douglas Shortcut}

For $u = x_1 x_2$ with equal exponents:

\begin{itemize}[leftmargin=*]
\item \textbf{Original:} Always spend half your income on each good
\begin{itemize}
\item $x_1^* = \frac{m}{2p_1}$, $x_2^* = \frac{m}{2p_2}$
\end{itemize}

\item \textbf{Compensated:} The $MRS = p_1'/p_2$ gives you $\frac{x_2}{x_1} = \frac{p_1'}{p_2}$
\begin{itemize}
\item Combine this with $x_1 x_2 = \bar{u}$ (the original utility)
\item Solve for $x_1$ and $x_2$
\end{itemize}

\item \textbf{New:} Spend half your income on each good again
\begin{itemize}
\item $x_1^{new} = \frac{m}{2p_1'}$, $x_2^{new} = \frac{m}{2p_2}$
\end{itemize}
\end{itemize}

\section*{Bottom Line}

When a price changes, your consumption changes for two reasons:

\begin{enumerate}
\item You substitute toward the relatively cheaper good (substitution effect)
\item You're effectively richer or poorer, so you adjust everything (income effect)
\end{enumerate}

The Slutsky decomposition just separates these two effects so you can see how much of the change came from each reason.

\vspace{0.5cm}

\begin{tcolorbox}[colback=yellow!20!white,colframe=orange!75!black]
\textbf{Remember:}
\begin{itemize}[leftmargin=*]
\item Three bundles: Original, Compensated, New
\item Compensated = "What if I kept you equally happy at new prices?"
\item Sub effect = Always opposes price change
\item Income effect = Depends on normal vs inferior
\item Total = Sub + Income (always!)
\end{itemize}
\end{tcolorbox}

\end{document}
